\documentclass[12pt]{article}
\usepackage{amsmath, amssymb, geometry}
\geometry{margin=1in}
\setlength{\parindent}{0pt}

\title{In-Class Exercise 1 Solutions: Single-Variable Optimization and Logarithms\\[7pt]
\large ECON 320 - Introduction to Mathematical Economics}
\author{Muhebullah Karimzada}
\date{Fall 2025}


\begin{document}
\maketitle

\vspace{1em}

\textbf{1. Maximization of Profit Function}

A firm's total revenue and cost functions are:
\[
R(Q) = 80Q - 2Q^2, \quad C(Q) = 20 + 10Q
\]

\textbf{Solution:}

\textit{(a) Profit function:}
\[
\pi(Q) = R(Q) - C(Q) = (80Q - 2Q^2) - (20 + 10Q) = -2Q^2 + 70Q - 20
\]

\textit{(b) First-order condition (FOC):}
\[
\frac{d\pi}{dQ} = -4Q + 70 = 0 \Rightarrow Q^* = 17.5
\]

\textit{(c) Maximum profit:}
\[
\pi(17.5) = -2(17.5)^2 + 70(17.5) - 20 = -612.5 + 1225 - 20 = 592.5
\]

\textit{(d) Second-order condition:}
\[
\frac{d^2\pi}{dQ^2} = -4 < 0 \Rightarrow \text{maximum holds.}
\]
\newpage
\vspace{1em}

\textbf{2. Utility Maximization Problem}

Given:
\[
U(x) = 4x - 0.5x^2
\]

\textbf{Solution:}

\textit{(a)} \( \frac{dU}{dx} = 4 - x = 0 \Rightarrow x^* = 4. \)

\textit{(b)} Second derivative: \( \frac{d^2U}{dx^2} = -1 < 0 \Rightarrow \) maximum.

\textit{(c)} \textit{Interpretation:} Utility increases with leisure up to 4 hours, then decreases due to diminishing satisfaction. Thus, \(x^* = 4\) is the optimal leisure time.

\vspace{1em}

\textbf{3. Logarithmic Utility and Diminishing Marginal Utility}

\[
U(x) = \ln(x)
\]

\textbf{Solution:}

\textit{(a)} \( U'(x) = \frac{1}{x}, \quad U''(x) = -\frac{1}{x^2} \)

\textit{(b)} \( U'(x) > 0 \) for \( x>0 \) and \( U''(x) < 0 \Rightarrow \) marginal utility is positive but diminishing.

\textit{(c)} Elasticity:
\[
E_{Ux} = \frac{U'(x) \cdot x}{U(x)} = \frac{1/x \cdot x}{\ln(x)} = \frac{1}{\ln(x)}
\]
At \( x=10 \), \( E_{Ux} = \frac{1}{\ln(10)} \approx 0.434. \)

\textit{Interpretation:} Utility grows less than proportionally with \(x\); elasticity below one implies diminishing satisfaction.

\vspace{1em}

\textbf{4. Production Function and Marginal Product}

\[
Q = 10L^{0.5}
\]

\textbf{Solution:}

\textit{(a)} \( MPL = \frac{dQ}{dL} = 5L^{-0.5} = \frac{5}{\sqrt{L}} \)

\textit{(b)} \( MPL \) declines as \( L \) increases, so diminishing marginal productivity begins immediately.

\textit{(c)} Profit: \( \pi = pQ - wL = 5(10L^{0.5}) - 2L = 50L^{0.5} - 2L \)

\textit{FOC:} \( \frac{d\pi}{dL} = 25L^{-0.5} - 2 = 0 \Rightarrow L^{-0.5} = \frac{2}{25} \Rightarrow \sqrt{L} = 12.5 \Rightarrow L^* = 156.25. \)

\textit{Maximum profit:} \( \pi(156.25) = 50(12.5) - 2(156.25) = 625 - 312.5 = 312.5. \)

\vspace{1em}

\textbf{5. Maximization with Logarithmic Revenue}

\[
P(Q) = 100 - 4Q, \quad C(Q) = 20Q + 10, \quad R(Q) = P(Q)\cdot Q
\]

\textbf{Solution:}

\textit{(a)} \( R(Q) = 100Q - 4Q^2 \)

\[
\Pi(Q) = \ln(100Q - 4Q^2) - (20Q + 10)
\]

\textit{(b)} FOC:
\[
\frac{d\Pi}{dQ} = \frac{100 - 8Q}{100Q - 4Q^2} - 20 = 0
\]

Simplify numerator:
\[
100 - 8Q = 20(100Q - 4Q^2)
\]
\[
100 - 8Q = 2000Q - 80Q^2
\]
\[
80Q^2 - 2008Q + 100 = 0
\]

Solve quadratic:
\[
Q^* = \frac{2008 \pm \sqrt{2008^2 - 4(80)(100)}}{160}
\]

\[
Q^* \approx \frac{2008 \pm 2006}{160}
\Rightarrow Q^* \approx 25.08 \text{ or } 0.0125
\]
The feasible value (ensuring \(R>0\)) is \( Q^* = 25.08. \)

\textit{(c)} \textit{Interpretation:} Taking logs places higher weight on proportional revenue changes, approximating risk-averse or growth-sensitive behavior.

\vspace{1em}

\textbf{6. Elasticity and Logarithmic Differentiation}

\[
y = 20x^{0.8}
\]

\textbf{Solution:}

Take logs:
\[
\ln y = \ln 20 + 0.8 \ln x
\]
Differentiate:
\[
\frac{1}{y}\frac{dy}{dx} = 0.8 \frac{1}{x} \Rightarrow \frac{dy}{dx} = 0.8\frac{y}{x}
\]
Elasticity:
\[
E_{yx} = \frac{dy}{dx}\cdot \frac{x}{y} = 0.8
\]
\textit{Interpretation:} A 1\% increase in \(x\) raises \(y\) by 0.8\%.

\vspace{1em}

\textbf{7. Application of Natural Logarithm Properties}

\[
\ln \left( \frac{x^3 y^2}{z^4} \right) = 3\ln x + 2\ln y - 4\ln z
\]

\textit{Interpretation:} Output elasticity w.r.t.\ inputs can be read directly from coefficients. If \(x\) doubles (\(+100\%\)), output increases by approximately \(3\ln 2 \approx 2.079\), holding \(y,z\) constant.

\vspace{1em}


\end{document}
